\documentclass[12pt]{article}

\usepackage[a4paper,margin=2.5cm]{geometry}

\usepackage[onehalfspacing]{setspace}

\usepackage{amsmath}
\usepackage{amssymb}
\usepackage{xcolor}
\usepackage{mathtools}
\usepackage{amsthm}
\newtheorem{definition}{Definition}
\newtheorem{example}{Example}

\setlength{\parindent}{0pt}

\begin{document}

\begin{titlepage}
    \centering

    \vspace*{\fill}

    {\Huge Basis and Dimension II}

    \vspace{1cm}

    {\Large Jin-Ren Ryan Wang\par}

    \vspace{0.5cm}

    {\large \today\par}

    \vspace*{\fill}
\end{titlepage}

\tableofcontents

\newpage

\section{Dimensions of standard spaces}

Recall that for a vector space $V$ over $F$,

\begin{itemize}
    \item $S$ is a basis for $V$ if $S$ is linearly independent such that
    $V=\operatorname{span}(S)$;

    \item $\dim_F(V)=|S|$.
\end{itemize}

(Sometimes we may suppress $\dim_F$ as just $\dim$ if the underlying field is clear.)

\subsection{Example}

\textbf{1. (Vectors)}

\[
V=F^n
\]

is a vector space over $F$.
\[
\dim_F(F^n)=n.
\]

The standard basis is
\[
\{\mathbf{e_1,e_2,\ldots,e_n}\},
\]

where

\[
\{\mathbf{e_1}=(1,0,\ldots,0),
\mathbf{e_2}=(0,1,\ldots,0),
\ldots,
\mathbf{e_n}=(0,0,\ldots,1)\}.
\]

\textbf{2. (Matrix spaces)}

\[
\dim_F(M_n(F))=n^2.
\]

The ``standard'' basis is
\[
\{E_{11},E_{12},\ldots,E_{nn}\},
\]

where $E_{ij}$ is the matrix whose $(i,j)$ entry is $1$ and all other entries are $0$.\\
i.e.

\[
E_{11}
=
\begin{pmatrix}
1&0&\cdots&0\\
0&0&\cdots&0\\
\vdots&\vdots&\ddots&\vdots\\
0&0&\cdots&0
\end{pmatrix},
\quad
E_{12}
=
\begin{pmatrix}
0&1&\cdots&0\\
0&0&\cdots&0\\
\vdots&\vdots&\ddots&\vdots\\
0&0&\cdots&0
\end{pmatrix},
\quad
\ldots
\quad
E_{nn}
=
\begin{pmatrix}
0&0&\cdots&0\\
0&0&\cdots&0\\
\vdots&\vdots&\ddots&\vdots\\
0&0&\cdots&1
\end{pmatrix}.
\]

\newpage

\textbf{3.}

Find a basis and the dimension of the following subspace of $\mathbb{R}^4$:

\[
W=
\left\{
(x,y,z,w)\in\mathbb{R}^4 :
x+y+z+w=0
\text{ and }
-x+2z+3w=0
\right\}.
\]

\textbf{Step 1: Simplifying the equations}

\[
\left\{
\begin{aligned}
x+y+z+w &= 0 \qquad (1)\\
-x+2z+3w &= 0 \qquad (2)
\end{aligned}
\right.
\]

Adding $(1)$ and $(2)$, we obtain

\[
\left\{
\begin{aligned}
x+y+z+w &= 0 \qquad (1)\\
y+3z+4w &= 0 \qquad (2)'
\end{aligned}
\right.
\]

\bigskip

\textbf{Step 2: Backward substitution}

By $(2)'$,
\[
y+3z+4w=0.
\]

Let
\[
z=s,\qquad w=t.
\]

Then
\[
y=-3s-4t.
\]

By $(1)$,
\[
x+y+z+w=0,
\]

so
\begin{align*}
x+(-3s-4t)+s+t &=0\\
x-2s-3t &=0\\
x &=2s+3t.
\end{align*}

Therefore
\[
W
=
\left\{
(2s+3t,-3s-4t,s,t) : s,t\in\mathbb R
\right\}.
\]

\[
=
\left\{
s(2,-3,1,0) + t(3,-4,0,1) : s,t\in\mathbb R
\right\}.
\]

\[
=
\operatorname{span}
\left\{
(2,-3,1,0),
(3,-4,0,1)
\right\}.
\]

As
\[
B=
\{(2,-3,1,0),(3,-4,0,1)\}
\]

spans $W$, and $B$ is linearly independent, by def, $B$ forms a basis of $W$.

Therefore,
\[
\dim_{\mathbb R}(W)=2.
\]

\textbf{4. (Polynomial of degree at most $n$)}

\[
\dim_F(F[x]_{\le n})=n+1.
\]

The ``standard'' basis is
\[
\{1,x,x^2,\ldots,x^n\}.
\]

\textbf{5. (Zero space)}

Recall that
\[
\{0\}
=
\operatorname{span}(\varnothing),
\]

where
\[
\varnothing
\]

denotes the empty set.

Therefore,
\[
\varnothing
\]

forms a basis of
\[
\{0\}.
\]

Hence,
\[
\dim_F(\{0\})
=
|\varnothing|
=
0.
\]

\textbf{6. (Complex numbers)}

\[
\mathbb C
=
\{
a\cdot 1+b\cdot\sqrt{-1} : a,b\in\mathbb R\}.
\]

Hence,
\[
\dim_{\mathbb R}(\mathbb C)=2.
\]

On the other hand,
\[
\mathbb C
=
\{
a\cdot 1
:
a\in\mathbb C
\}.
\]

Hence,
\[
\dim_{\mathbb C}(\mathbb C)=1.
\]

\newpage

\section{Consequence of Replacement Theorem}

\subsection{Theorem}

\textbf{Theorem 3.2.1.}

Let $V$ be a vector space over $F$ of dimension $n$, and let
$S\subseteq V$.

\begin{enumerate}
    \item If $|S|>n$, then $S$ cannot be linearly independent.

    \item If $|S|<n$, then $S$ cannot span $V$.
\end{enumerate}

(1)
\begin{proof}

Given $|S|>n$.

Suppose, for contradiction, that $S$ is linearly independent.

Let $B$ be a basis of $V$. Then
\[
|B|=n.
\]

Since $B$ spans $V$, the Replacement Theorem implies
\[
|S|\le |B|=n,
\]

which contradicts $|S|>n$.

Therefore, $S$ cannot be linearly independent.

\end{proof}

(2)
\begin{proof}
Given $|S|<n$.

Suppose, for contradiction, that $S$ spans $V$.

Let $B$ be a basis of $V$. Then
\[
|B|=n.
\]

Since $B$ is linearly independent and $S$ spans $V$, the Replacement Theorem implies
\[
|B|\le |S|.
\]

Hence
\[
n=|B|\le |S|,
\]

which contradicts $|S|<n$.

Therefore, $S$ cannot span $V$
\end{proof}

\newpage

\textbf{Theorem 3.2.2.}
Let $V$ and $W$ be finite-dimensional vector spaces.

If
\begin{enumerate}
    \item $W\subseteq V$,
    \item $\dim_F(W)=\dim_F(V)$,
\end{enumerate}

then
\[
W=V.
\]

\begin{proof}
Given
\[
W\subseteq V
\]

and
\[
\dim_F(W)=\dim_F(V)=n.
\]

We WTS that $V\subseteq W$.

Suppose not. Then there exists
\[
\mathbf{v}\in V
\]

such that
\[
\mathbf{v}\notin W.
\]

Let
\[
S=\{\mathbf{w_1,\ldots,w_n}\}
\]

be a basis of $W$.

Since
\[
W=\operatorname{span}(S),
\]

we have
\[
\mathbf{v}\notin\operatorname{span}(S).
\]

Therefore,
\[
S'=\{\mathbf{w_1,\ldots,w_n,v}\}
\]

is linearly independent in $V$.

Since
\[
\dim_F(V)=n,
\]

Theorem 3.2.1 implies that any linearly independent subset\\
of $V$ contains at most $n$ vectors.

However,
\[
|S'|=n+1,
\]

a contradiction.

Hence
\[
V\subseteq W.
\]

Together with
\[
W\subseteq V,
\]

we conclude that
\[
W=V.
\]
\end{proof}

\textbf{Theorem 3.2.3.}
Let $V$ be a vector space of dimension $n$. Then any $n$ linearly independent vectors form a basis of $V$.

(In literature, we say that any maximal linearly independent subset forms a basis.)

\begin{proof}
Let $S$ be a linearly independent subset of $V$ with
\[
|S|=n.
\]

We WTS that
\[
\operatorname{span}(S)=V.
\]

Let
\[
W=\operatorname{span}(S).
\]

Then $W$ is a subspace of $V$, so
\[
W\subseteq V.
\]

Since $S$ spans $W$ and is linearly independent, $S$ is a basis of $W$.

Therefore,
\[
\dim_F(W)=|S|=n.
\]

Since
\[
\dim_F(V)=n,
\]

we have
\[
\dim_F(W)=\dim_F(V).
\]

By Theorem 3.2.2,
\[
W=V.
\]

Hence
\[
\operatorname{span}(S)=V.
\]

Therefore $S$ is a basis of $V$.

\end{proof}

\subsection{Example}

\textbf{1.} 
$S_1=\{(1,0,1),(1,1,0),(2,3,1),(-1,2,3)\}$ cannot be linearly independent.

\[
S_1=\{(1,0,1),(1,1,0),(2,3,1),(-1,2,3)\}.
\]

Since
\[
\dim_{\mathbb R}(\mathbb R^3)=3,
\]

Theorem 3.2.1 implies that any linearly independent subset of $\mathbb R^3$ contains at most $3$ vectors.

However,
\[
|S_1|=4.
\]

Therefore $S_1$ cannot be linearly independent.

\bigskip

\textbf{2.}
$S_2=\{(1,0,1),(1,1,0)\}$ cannot span $\mathbb{R}^3$.

\[
S_2=\{(1,0,1),(1,1,0)\}.
\]

Since
\[
\dim_{\mathbb R}(\mathbb R^3)=3,
\]

Theorem 3.2.1 implies that any spanning subset of$\mathbb R^3$ must contain at least $3$ vectors.

However,
\[
|S_2|=2.
\]

Therefore $S_2$ cannot span $\mathbb R^3$.

\bigskip

\textbf{3.}
The subset $S=\{(1,0,1),(1,1,0),(2,3,1)\}$ forms a basis of $\mathbb{R}^3$.

(Check.) $S$ is linearly independent.

Since
\[
\dim_{\mathbb R}(\mathbb R^3)=3=|S|,
\]

$S$ is a \textcolor{red}{maximal linearly independent set}. Therefore $S$ is a basis of $\mathbb R^3$.

Remark: This avoids the tedious check of spanning.

\bigskip

\textbf{4.}
Let $V=M_2(\mathbb{R})$ and consider the subset
\[
S=
\left\{
A_1=
\begin{pmatrix}
1&1\\
0&0
\end{pmatrix},
\;
A_2=
\begin{pmatrix}
1&0\\
1&0
\end{pmatrix},
\;
A_3=
\begin{pmatrix}
1&0\\
0&1
\end{pmatrix},
\;
A_4=
\begin{pmatrix}
0&1\\
0&1
\end{pmatrix},
\;
A_5=
\begin{pmatrix}
0&1\\
0&1
\end{pmatrix},
\;
A_6=
\begin{pmatrix}
0&0\\
1&1
\end{pmatrix}
\right\}.
\]

(a)
Prove that $\operatorname{span}(S) = V$

\newpage

\begin{proof}
Let
\[
S'=\{A_1,A_2,A_3,A_4\}.
\]

Suppose
\[
a_1A_1+a_2A_2+a_3A_3+a_4A_4
=
\begin{pmatrix}
0&0\\
0&0
\end{pmatrix}.
\]

By comparing the entries, we obtain
\[
\left\{
\begin{aligned}
a_1+a_2+a_3 &=0,\\
a_1+a_4 &=0,\\
a_2 &=0,\\
a_3+a_4 &=0.
\end{aligned}
\right.
\]

Hence
\[
a_1=a_2=a_3=a_4=0.
\]

Therefore $S'$ is linearly independent.

Since
\[
|S'|=4
=
\dim_{\mathbb R}(M_2(\mathbb R)),
\]

Theorem 3.2.3 implies that $S'$ is a basis of $V$.

Thus
\[
V=\operatorname{span}(S').
\]

Since
\[
S'\subseteq S,
\]

we conclude that
\[
\operatorname{span}(S)=V.
\]
\end{proof}

(b)
Find a subset of $S$ that is a basis of $V$

By part (a),
\[
\{A_1,A_2,A_3,A_4\}
\]
is a basis of $V$.\\

(c)
Find a subset of $S$ with 4 elements which is not a basis of $V$

Observe that
\[
A_1+A_6=A_2+A_4
=
\textcolor{red}{
\begin{pmatrix}
1&1\\
1&1
\end{pmatrix}.
}
\]

Therefore
\[
A_1-A_2-A_4+A_6
=
\begin{pmatrix}
0&0\\
0&0
\end{pmatrix}.
\]

Since the coefficients are not all zero,
\[
\{A_1,A_2,A_4,A_6\}
\]

is linearly dependent.

Hence
\[
\{A_1,A_2,A_4,A_6\}
\]

is not a basis of $V$.

\newpage

\section{Sifiting and proof of Replacement Theorem}

\subsection{Theorem}

\textbf{Theorem 3.3.1.}
A subset
\[
S=\{\mathbf{v_1,v_2,\ldots,v_m}\}
\]
of $V$ is linearly independent if and only if for each
$r=1,\ldots,m$,
$v_r$ is not a linear combination of
$\mathbf{v_1,v_2,\ldots,v_{r-1}}$.
In particular, this requires
$\mathbf{v_1}\neq\mathbf{0}$.

\begin{proof}
This is just Theorem 2.2.1.
\end{proof}

\textbf{Theorem 3.3.2.}
Let
\[
S=\{\mathbf{v_1,v_2,\ldots,v_m,w}\}.
\]

If $w$ is a linear combination of
$\mathbf{v_1,v_2,\ldots,v_m}$, then
\[
\operatorname{span}(S)=\operatorname{span}\{\mathbf{v_1,v_2,\ldots,v_m}\}.
\]

\begin{proof}
\textcolor{red}{$(\subseteq)$}\\
Given
\[
\mathbf{w}=a_1\mathbf{v_1}+a_2\mathbf{v_2}+\cdots+a_m\mathbf{v_m}
\]

for some $a_i\in F$.

We first show
\[
\operatorname{span}(S)\subseteq\operatorname{span}\{\mathbf{v_1,v_2,\ldots,v_m}\}.
\]

Let $\mathbf{v}\in\operatorname{span}(S)$. Then

\begin{align*}
\mathbf{v}
&=b_1\mathbf{v_1}+b_2\mathbf{v_2}+\cdots+b_m\mathbf{v_m}+c\mathbf{w}\\
&=b_1\mathbf{v_1}+b_2\mathbf{v_2}+\cdots+b_m\mathbf{v_m}
+c(a_1\mathbf{v_1}+\cdots+a_m\mathbf{v_m})\\
&=(b_1+ca_1)\mathbf{v_1}+(b_2+ca_2)\mathbf{v_2}+\cdots+(b_m+ca_m)\mathbf{v_m}.
\end{align*}

Hence
\[
\mathbf{v}\in\operatorname{span}\{\mathbf{v_1,v_2,\ldots,v_m}\},
\]

so
\[
\operatorname{span}(S)\subseteq\operatorname{span}\{\mathbf{v_1,v_2,\ldots,v_m}\}.
\]

\textcolor{red}{$(\subseteq)$}\\
Conversely, let
\[
\mathbf{v}\in\operatorname{span}\{\mathbf{v_1,v_2,\ldots,v_m}\}.
\]

Then
\begin{align*}
\mathbf{v}
&=d_1\mathbf{v_1}+d_2\mathbf{v_2}+\cdots+d_m\mathbf{v_m}\\
&=\underbrace{d_1\mathbf{v_1}+d_2\mathbf{v_2}+\cdots+d_m\mathbf{v_m}+0\cdot \mathbf{w}
}_{\text{linear combination of $S$}}.
\end{align*}

Hence
\[
\mathbf{v}\in\operatorname{span}(S),
\]

which implies
\[
\operatorname{span}\{\mathbf{v_1,v_2,\ldots,v_m}\}\subseteq\operatorname{span}(S).
\]

Therefore,
\[
\operatorname{span}(S)=\operatorname{span}\{\mathbf{v_1,v_2,\ldots,v_m}\}.
\]
\end{proof}

\textbf{Theorem 3.3.3. (Sifiting method)}
Let
\[
S=\{\mathbf{v_1,v_2,\ldots,v_m}\}
\]

be a subset of $V$. Then there exists a subset
\[
S'\subseteq S
\]

such that

\begin{enumerate}
    \item $S'$ is linearly independent; and
    \item $\operatorname{span}(S')=\operatorname{span}(S).$
\end{enumerate}

\begin{proof}
For each $r$, check whether $\mathbf{v_r}$ is a linear combination of the
previous vectors
\[
\mathbf{v_1,v_2,\ldots,v_{r-1}}.
\]

\begin{enumerate}
    \item Keep $\mathbf{v_r}$ in $S$ if $\mathbf{v_r}$ is \emph{not} a linear combination of
    $\mathbf{v_1,v_2,\ldots,v_{r-1}}$ (not redundant).

    \item Remove $\mathbf{v_r}$ from $S$ if $\mathbf{v_r}$ \emph{is} a linear combination of
    $\mathbf{v_1,v_2,\ldots,v_{r-1}}$ (redundant).
\end{enumerate}

After this process, we obtain a subset $S'$ in which no element is a
linear combination of the preceding members. Therefore, by Theorem~3.3.1, $S'$ is linearly independent.
Moreover, by Theorem~3.3.2, removing redundant members from $S$ does not
change its linear span. Hence, we have $\operatorname{span}(S')=\operatorname{span}(S)$.

\end{proof}

\textbf{Theorem 3.3.4. (Basis extension theorem)}

Let $V$ be a finite-dimensional vector space, and let $S$ be a linearly independent subset of $V$. Then $S$ can be extended into a basis of $V$.

\begin{proof}
Let
\[
S=\{\mathbf{v_1,\ldots,v_m}\}
\]

be a linearly independent subset of $V$.

Let
\[
B=\{b_1,\ldots,b_n\}
\]

be any basis of $V$.

Consider
\[
S\cup B = \{\mathbf{v_1,\ldots,v_m,b_1,\ldots,b_n}\}.
\]

Apply the Sifting Method to $S\cup B$.

Since $S$ is linearly independent, none of the vectors
$\mathbf{v_1,\ldots,v_m}$ will be removed during the sifting process.

Hence we obtain
\[
S\cup B',\qquad\text{where }B'\subseteq B.
\]

Notice that

\begin{enumerate}
    \item $S\cup B'$ is linearly independent;
    \item $\operatorname{span}(S\cup B')= \operatorname{span}(S\cup B).$
\end{enumerate}

Since $B$ is a basis of $V$,
\[
\operatorname{span}(S\cup B)=V.
\]

Therefore,
\[
\operatorname{span}(S\cup B')=V.
\]

Hence $S\cup B'$ is both linearly independent and spans $V$.
Therefore,
\[
S\cup B'
\]

forms a basis of $V$, which is an extension of $S$.
\end{proof}

\newpage

\subsection{Proof of Exchange/Replacement Theorem}
This is the proof of Replacement Theorem (Theorem 2.5.1) which says that if
\begin{itemize}
    \item
    $S_1=\{\mathbf{v}_1,\mathbf{v}_2,\ldots,\mathbf{v}_n\}$ is a spanning set of $V$.

    \item
    $S_2=\{\mathbf{w}_1,\mathbf{w}_2,\ldots,\mathbf{w}_m\}$ is linearly independent.
\end{itemize}

Then
\[
m\le n.
\]

\begin{proof}
Idea: Insert $\mathbf{w}_j$ one-by-one into $S_1$, then apply the Sifting Method.

Consider
\[
B_1
=
\{\textcolor{red}{\mathbf{w}_1},\mathbf{v}_1,\mathbf{v}_2,\ldots,\mathbf{v}_n\}
=
\{\textcolor{red}{\mathbf{w}_1}\}\cup S_1.
\]

Applying the Sifting Method to $B_1$, we must remove some $\mathbf{v}_i$'s. Hence we obtain
\[
B_1' =\{\textcolor{red}{\mathbf{w}_1}\}\cup S_1', \text{ where } |S_1'|<|S_1|.
\]

Moreover,
\[
B_1'
\]
remains a spanning set of $V$.

Now consider
\[
B_2 = \{\textcolor{red}{\mathbf{w}_2,\mathbf{w}_1}\}\cup S_1'.
\]

Applying the Sifting Method again, none of the $\mathbf{w}_i$'s will be removed because
\[
S_2 = \{\mathbf{w}_1,\ldots,\mathbf{w}_m\}
\]
is linearly independent.

Hence we obtain
\[
B_2' = \{\textcolor{red}{\mathbf{w}_2,\mathbf{w}_1}\}\cup S_1'', \text{ where } |S_1''|<|S_1'|.
\]

Continue this process. Finally we obtain
\[
B_m' = \{\textcolor{red}{\mathbf{w}_m,\mathbf{w}_{m-1},\ldots,\mathbf{w}_1}\} \cup S_1^{(m)},
\]
where
\[
0 \le |S_1^{(m)}| < |S_1^{(m-1)}| < \cdots < |S_1''| < |S_1'| < |S_1| = n.
\]

Thus there are at least $\textcolor{red}{m}$ distinct integers strictly less than $n$. 

Therefore, we must have
\[
\textcolor{red}{m}\le n.
\]
\end{proof}

\section{Direct sum of subspaces}

\subsection{Definition}

\textbf{Definition 3.4.1. (Sum of subspaces)}

Let $W_1$, $W_2$ be subspaces of a vector space $V$. Then define
\begin{align*}
    W_1 + W_2 \coloneq \{\mathbf{w_1} + \mathbf{w_2} : \mathbf{w_1}\in W_1 \text{ and } \mathbf{w_2}\in W_2\} 
\end{align*}

\textbf{Definition 3.4.2. (Direct Sum)}

We write $W_1 + W_2$ = $W_1 \bigoplus W_2$ if $W_1 \cap W_2$ = \{$\mathbf{0}$\}. In this case, \\
we call it the \emph{direct sum} (instead of just sum) of $W_1$ and $W_2$.

\subsection{Theorem}

\textbf{Theorem 3.4.1.}

Let $W_1 + W_2$ be a subspaces of $V$. Indeed, this is smallest subspace\\ 
containing both $W_1$ and $W_2$

\begin{proof}
We verify the three subspace conditions.

\bigskip

\fbox{Non-empty}

Since $W_1$ and $W_2$ are subspaces,
\[
\textcolor{red}{\mathbf{0}}\in \textcolor{red}{W_1},\qquad \textcolor{green}{\mathbf{0}}\in \textcolor{green}{W_2}.
\]
Hence
\[
\mathbf{0}=\textcolor{red}{\mathbf{0}}+\textcolor{green}{\mathbf{0}}\in \textcolor{red}{W_1}+\textcolor{green}{W_2}. \implies W_1 + W_2 \text{ is non-empty.}
\]

\fbox{Closed under addition}

Let
\[
\mathbf{a},\mathbf{b}\in W_1+W_2.
\]
Then
\[
\left\{
\begin{aligned}
\mathbf{a}&=\mathbf{w}_1+\mathbf{w}_2 \quad \text{ for some } \mathbf{w}_1\in W_1 \text{ and } \mathbf{w}_2\in W_2. \\
\mathbf{b}&=\textcolor{red}{\mathbf{w}_1'+\mathbf{w}_2'} \quad \text{ for some } \textcolor{red}{\mathbf{w}_1'}\in W_1 \text{ and } \textcolor{red}{\mathbf{w}_2'}\in W_2.
\end{aligned}
\right.
\]

Therefore,
\[
\begin{aligned}
\mathbf{a}+\mathbf{b}
&=(\mathbf{w}_1+\textcolor{red}{\mathbf{w}_1'})+(\mathbf{w}_2+\textcolor{red}{\mathbf{w}_2'}).
\end{aligned}
\]

Since $W_1$ and $W_2$ are subspaces,
\[
\mathbf{w}_1+\textcolor{red}{\mathbf{w}_1'}\in W_1, \qquad \mathbf{w}_2+\textcolor{red}{\mathbf{w}_2'}\in W_2.
\]
Hence
\[
\mathbf{a}+\mathbf{b}\in W_1+W_2. \implies W_1 + W_2 \text{ is Closed under addition}
\]

\fbox{Closed under scalar multiplication}

Let
\[
\mathbf{a}\in W_1+W_2, \text{ and } c \text{ is a scalar}
\]

Then
\[
\left\{
\begin{aligned}
\mathbf{a}&=\mathbf{w}_1+\mathbf{w}_2 \text{ for some } \mathbf{w_1}\in W_1 \text{ and } \mathbf{w_2}\in W_2\\
c\mathbf{a}&=c(\mathbf{w}_1+\mathbf{w}_2)=c\mathbf{w_1}+c\mathbf{w_2} \text{ for some } \mathbf{w_1}\in W_1 \text{ and } \mathbf{w_2}\in W_2        
\end{aligned}
\right.
\]

Since
\[
c\mathbf{w}_1\in W_1,\qquad c\mathbf{w}_2\in W_2,
\]
we conclude that
\[
c\mathbf{a}\in W_1+W_2.
\]

Therefore, $W_1+W_2$ is a subspace of $V$.
\end{proof}

\textbf{Theorem 3.4.2.}

Let $V$ be a finite-dimensional vector space, and let $W_1,W_2$ be subspaces of $V$. Then
\[
\dim(W_1+W_2)=\dim(W_1)+\dim(W_2)-\dim(W_1\cap W_2).
\]

\begin{proof}

\fbox{Idea. Skill for dimension formula: Start with the ``smallest'' subspace}

\[
\fcolorbox{red}{white}{\textcolor{red}{$W_1\cap W_2\subseteq W_1,\;W_2 \subseteq W_1+W_2.$}}
\]

Let
\[
n=\dim(W_1\cap W_2),
\]
and let
\[
B=\{\mathbf{b}_1,\mathbf{b}_2,\ldots,\mathbf{b}_n\}
\]
be a basis of $W_1\cap W_2$.

Since $B\subseteq \textcolor{red}{W_1}$ and $B$ is linearly independent,
by the Basis Extension Theorem, we can extend $B$ to a basis of $W_1$:
\[
B_1= \{\mathbf{b}_1,\ldots,\mathbf{b}_n, \textcolor{red}{\mathbf{s}_1,\ldots,\mathbf{s}_p}\}, \text{ a basis of } \textcolor{red}{W_1}
\]

Similarly, since $B\subseteq \textcolor{red}{W_2}$ and $B$ is linearly independent,
we may extend $B$ to a basis of $W_2$:
\[
B_2= \{\mathbf{b}_1,\ldots,\mathbf{b}_n, \mathbf{t}_1,\ldots,\mathbf{t}_q\}, \text{ a basis of } \textcolor{red}{W_2}
\]

\fbox{BIG CLAIM}

Let
\[
S= \{\mathbf{b}_1,\ldots,\mathbf{b}_n, \textcolor{red}{\mathbf{s}_1,\ldots,\mathbf{s}_p}, \mathbf{t}_1,\ldots,\mathbf{t}_q\}.
\]

We claim that $S$ forms a basis of $W_1+W_2$.

\medskip

\begin{proof}[Proof of the claim]
We verify the three subspace conditions:

(1) \fbox{$S$ spans $W_1+W_2$}

Let
\[
\mathbf{w}\in W_1+W_2.
\]

Then there $\exists$ $\mathbf{w}_1\in W_1$ and $\mathbf{w}_2\in W_2$ s.t.

\[
\mathbf{w}=\mathbf{w}_1+\mathbf{w}_2 \in W_1+W_2
\]

Since $B_1$ is a basis of $W_1$,
\[
\mathbf{w}_1= \sum_{i=1}^{n}\alpha_i\mathbf{b}_i + \sum_{j=1}^{p}\beta_j\textcolor{red}{\mathbf{s}_j}.
\]

And $B_2$ is a basis of $W_2$,
\[
\mathbf{w}_2= \sum_{i=1}^{n}\alpha_i'\mathbf{b}_i + \sum_{k=1}^{q}\gamma_k\mathbf{t}_k.
\]

So
\[
\begin{aligned}
\mathbf{w}
&=
\mathbf{w}_1+\mathbf{w}_2\\
&=
\sum_{i=1}^{n}(\alpha_i+\alpha_i')\mathbf{b}_i+\sum_{j=1}^{p} \beta_j\textcolor{red}{\mathbf{s}_j} + \sum_{k=1}^{q}\gamma_k\mathbf{t}_k.
\end{aligned}
\]

Therefore,
\[
\mathbf{w}\in\operatorname{span}(S).
\]

Hence $S$ spans $W_1+W_2$.

\medskip

(2) \fbox{$S$ is linearly independent.}

Suppose $\exists$ scalars $\alpha_i$, \textcolor{red}{$\beta_j$}, $\gamma_k$ s.t.
\[
\sum_{i=1}^{n}\alpha_i\mathbf{b}_i + \sum_{j=1}^{p}\textcolor{red}{\beta_j\mathbf{s}_j} + \sum_{i=1}^{q}\gamma_k\mathbf{t}_k= \mathbf0.
\]

Rearranging,
\[
\sum_{i=1}^{n}\alpha_i\mathbf{b}_i + \sum_{j=1}^{p}\textcolor{red}{\beta_j\mathbf{s}_j} = -\sum_{k=1}^{q}\gamma_k\mathbf{t}_k.
\]

Let $\mathbf{v}$ be left-hand side. Undoubtedly, $\mathbf{v}$ equals right-hand side.

Clearly, $\mathbf{v}$ belongs to both $W_1$ and $W_2$. Hence both sides belong to $W_1\cap W_2$.

As
\[
B=\{\mathbf{b}_1,\ldots,\mathbf{b}_n\}
\]
is a basis of $W_1\cap W_2$, $\mathbf{v}$ can be written as
\[
\mathbf{v} = \sum_{i=1}^{n}\sigma_i\mathbf{b}_i \implies \mathbf{0} =  \sum_{i=1}^{n}\sigma_i\mathbf{b}_i + \sum_{k=1}^{q}\gamma_k\mathbf{t_k}
\]

As $B_2$ = \{$\mathbf{b}_1, \ldots \mathbf{b}_n,$, $\mathbf{t}_1, \ldots \mathbf{t}_q$\} is linearly independent,
this forces \fcolorbox{red}{white}{$\sigma_i$ = 0 $\forall i$}, \fcolorbox{red}{white}{$\gamma_k$ = 0 $\forall k$}

Then the rearranging part above can be reduced to
\[
\sum_{i=1}^{n}\alpha_i\mathbf{b}_i + \sum_{j=1}^{p}\textcolor{red}{\beta_j\mathbf{s}_j} = \mathbf{0}.
\]

As $B_1$ = \{$\mathbf{b}_1, \ldots \mathbf{b}_n,$, $\mathbf{s}_1, \ldots \mathbf{s}_p$\} is linearly independent,
this forces \fcolorbox{red}{white}{$\alpha_i$ = 0 $\forall i$}, \fcolorbox{red}{white}{$\beta_j$ = 0 $\forall j$}

So $S$ is linearly independent.

\end{proof}

\bigskip

Hence $S$ is a basis of $W_1+W_2$,
\[
\begin{aligned}
\dim(W_1+W_2)
&=
|S|\\
&=
n+p+q\\
&=
(n+p)+(n+q)-n\\
&=
\dim(W_1)+\dim(W_2)-\dim(W_1\cap W_2).
\end{aligned}
\]

\end{proof}

\textbf{Theorem 3.4.3.}

If $W_1,W_2$ are finite-dimensional subspaces of $V$ with $W_1\cap W_2=\{\mathbf{0}\}$,

then
\[
\dim(W_1\oplus W_2)
=
\dim(W_1)+\dim(W_2).
\]

\begin{proof}

Since
\[
W_1\oplus W_2=W_1+W_2,
\]
By Theorem 3.4.2, we have
\[
\begin{aligned}
\dim(W_1\oplus W_2)
&=\dim(W_1+W_2) \\
&=\dim(W_1)+\dim(W_2)-\dim(W_1\cap W_2) \\
&=\dim(W_1)+\dim(W_2)-\textcolor{red}{0} \\
&=\dim(W_1)+\dim(W_2).
\end{aligned}
\]

\end{proof}

\newpage

\textbf{Theorem 3.4.4.}

Let $V$ be a finite-dimensional vector space and let $W$ be a subspace of $V$.
Then there exists another subspace $W'$ of $V$ such that $V=W\oplus W'$.\\

In this case, we call $W'$ A \emph{complementary subspace} of $W$ in $V$.

\begin{proof}

Let
\[
S=\{\mathbf{v}_1,\mathbf{v}_2,\ldots,\mathbf{v}_k\}
\]

be a basis of $W$.

By the Basis Extension Theorem, we may extend \textcolor{red}{$S$} to a basis of $V$, and we say that
\[
\textcolor{red}{S'}= \{\underbrace{\mathbf{v}_1,\ldots,\mathbf{v}_k}_{\textcolor{red}{\text{span W}}}, \mathbf{s}_1,\ldots,\mathbf{s}_n \}.
\]

Define
\[
W'= \operatorname{span} \{\mathbf{s}_1,\ldots,\mathbf{s}_n \}.
\]

Let $V$ = $W \bigoplus W'$, we claim that
\[
\text{(1) }V\subseteq W+W',
\qquad
\text{(2) }W\cap W'=\{\mathbf0\}.
\]

\begin{proof}[proof of (1)]
Let
\[
\mathbf{v}\in V.
\]

As $\textcolor{red}{S'}$ is a basis of $V$,
\[
\mathbf{v} = \underbrace{\sum_{i=1}^{k}\alpha_i\mathbf{v}_i}_{\in W} + \underbrace{\sum_{j=1}^{n}\beta_j\mathbf{s}_j}_{\in W'} \in W + W'
\]

\end{proof}

\begin{proof}[proof of (2)]
Let
\[
\mathbf{w}\in W\cap W'.
\]

Since both $\mathbf{w}\in W$ and $\mathbf{w}\in W'$,

\[
\mathbf{w} = \sum_{i=1}^{k}\alpha_i\mathbf{v}_i = \sum_{j=1}^{n}\beta_j\mathbf{s}_j.
\]

Hence,
\[
\sum_{i=1}^{k}\alpha_i\mathbf{v}_i-\sum_{j=1}^{n}\beta_j\mathbf{s}_j=\mathbf0.
\]

Since \textcolor{red}{$S'$} is linearly independent,
\[
\alpha_i=\beta_j=0, \quad \forall i,j
\]

Thus, $\mathbf{w}=\mathbf0,$  $\implies$ $W\cap W'=\{\mathbf0\}$.

\end{proof}

\end{proof}

\subsection{Question}
Question: Is true that, $\dim(W_1+W_2+W_3) = \dim(W_1)+\dim(W_2)+\dim(W_3)-\dim(W_1\cap W_2)-\dim(W_2\cap W_3)-\dim(W_3\cap W_1)+\dim(W_1\cap W_2\cap W_3)?$

Answer: \textcolor{red}{False}.

\subsection{Example}

\subsubsection{Direct sum of symmetric and skew-symmetric matrices}
Consider the vector space $V=M_n(\mathbb{C})$ and two of its subspaces :
\begin{itemize}
    \item $W_1=\{A\in M_n(\mathbb{C})\mid A^t=A\}$
    \item $W_2=\{A\in M_n(\mathbb{C})\mid A^t=-A\}$
\end{itemize}

Prove that $V=W_1\oplus W_2$.

\begin{proof}
We need to show
\begin{itemize}
    \item[(1)]$V=W_1+W_2$\\
    i.e. any square matrix is a sum of a symmetric and a skew-symmetric matrix
    \item[(2)]$W_1\cap W_2=\{\textcolor{red}{\mathbf{0}}\} \quad$ ($\textcolor{red}{\text{zero matrix}}$) 
\end{itemize}

\medskip

\begin{proof}[proof of (1)]
Let
\[
A\in M_n(\mathbb C).
\]

Observe that
\[
A=\textcolor{red}{\frac12(A+A^t)}+\frac12(A-A^t).
\]

Let
\[
B=\frac12(A+A^t), \qquad C=\frac12(A-A^t).
\]

Then
\[
B^t = \frac12(A+A^t)^t = \frac12(A^t+A) = B,
\]

so $B$ is symmetric.

\medskip

Similarly,
\[
C^t = \frac12(A-A^t)^t = \frac12(A^t-A) = -C,
\]

so $C$ is skew-symmetric

\medskip

Therefore,
\[
A=B+C\in W_1+W_2,
\]

which implies
\[
V=W_1+W_2.
\]

\end{proof}

\medskip

\begin{proof}[proof of (2)]
Let
\[
A\in W_1\cap W_2.
\]

Then
\[
A^t=A,\quad A^t=-A.
\]

Hence,
\[
A=-A,
\]

so
\[
2A=\textcolor{red}{\mathbf{0}}.
\]

Since the underlying field is $\mathbb C$,
\[
A=\textcolor{red}{\mathbf{0}}.
\]

Therefore,
\[
W_1\cap W_2=\{\textcolor{red}{\mathbf{0}}\}.
\]

\end{proof}

\medskip

Combining (1) and (2), we conclude that
\[
V=W_1 \textcolor{red}{\oplus} W_2.
\]

\end{proof}

\end{document}